\documentclass[12pt,a4paper]{article}

% --- Encodage et langue ---
\usepackage[utf8]{inputenc}
\usepackage[T1]{fontenc}
\usepackage{lmodern}

% --- Maths ---
\usepackage{amsmath,amssymb,amsthm}

% --- Mise en page ---
\usepackage[top=2.5cm,bottom=2.5cm,left=2.8cm,right=2.8cm]{geometry}
\usepackage{parskip}

% --- Tableaux ---
\usepackage{booktabs}
\usepackage{array}
\usepackage{multirow}

% --- Figures et couleurs ---
\usepackage{tikz}
\usetikzlibrary{arrows.meta, positioning, decorations.pathreplacing, calc, fit, backgrounds}
\usepackage{xcolor}

% --- Liens ---
\usepackage{hyperref}
\hypersetup{colorlinks=true, linkcolor=blue!60!black, citecolor=blue!60!black}

% --- Environnements théorème ---
\newtheorem{theorem}{Théorème}[section]
\newtheorem{lemma}[theorem]{Lemme}
\newtheorem{corollary}[theorem]{Corollaire}
\theoremstyle{remark}
\newtheorem{remark}{Remarque}[section]
\newtheorem{example}{Exemple}[section]

% --- Macros ---
\newcommand{\Inc}{\ensuremath{\mathrm{Include}}}
\newcommand{\Exc}{\ensuremath{\mathrm{Exclude}}}
\newcommand{\vbar}{\bar{v}}
\newcommand{\ps}{\text{ p.s.}}
\newcommand{\DF}{D_F}

% ======================================================
\begin{document}
% ======================================================

\begin{center}
  {\LARGE\bfseries Convergence de la clause conditionnée}\\[6pt]
  {\large Généralisation multi-features de la règle à 4 cas}\\[4pt]
  {\large\itshape Littéral pertinent, neutralité des littéraux non pertinents,
  et fenêtre admissible pour $S$}\\[12pt]
  {\normalsize Analyse par chaînes de Markov --- suite de
  \textit{Démonstration de convergence de la règle d'apprentissage à 4 cas}}
\end{center}

\vspace{1em}
\hrule
\vspace{1.5em}

\tableofcontents

\newpage

% ======================================================
\section{Objectif et portée}
% ======================================================

Le document \textit{Démonstration de convergence de la règle d'apprentissage à 4 cas}
établit la convergence presque sûre (cas sans bruit) et la concentration de la
distribution stationnaire (cas bruité) pour \textbf{une seule feature} $x\in\{0,1\}$,
gouvernée par une paire d'automates $(v,\vbar)$.

Notre architecture complète conditionne chaque clause sur $K$ features tirées au
hasard et figées, et laisse les $n-K$ features restantes libres, apprises par la
règle à 4 cas sur le sous-ensemble d'exemples où les $K$ conditions sont satisfaites.
Ce document répond à la question laissée ouverte : \textbf{que se passe-t-il quand
plusieurs automates libres apprennent simultanément dans une même clause ?}

\begin{center}
\renewcommand{\arraystretch}{1.3}
\begin{tabular}{p{6.3cm}p{6.3cm}}
\toprule
\textbf{Ce que ce document prouve} & \textbf{Ce que ce document ne prouve pas} \\
\midrule
Une fois que le conditionnement à $K$ features a isolé un unique littéral libre
pertinent, la clause converge vers l'implémentation correcte de ce littéral, et
tous les autres littéraux libres se neutralisent (deviennent tautologiques).
&
Que le tirage aléatoire des $K$ features conditionnées trouve effectivement un tel
littéral pertinent --- c'est un phénomène combinatoire (la « loterie »),
caractérisé empiriquement (Mux11, Iris, Digits, MNIST), pas démontré ici.
\\
\addlinespace
Une fenêtre explicite $S\in(S_{\text{bas}}, S_{\max})$ garantissant cette double
convergence, avec vitesse géométrique explicite.
&
La convergence directe vers une cible à $m\ge2$ littéraux \emph{sans}
conditionnement --- empiriquement hors de portée de la règle à 4 cas asymétrique
(échec de \texttt{hybrid\_coord\_4case.cpp}), c'est le régime que couvre PCL, pas nous.
\\
\bottomrule
\end{tabular}
\end{center}

% ======================================================
\section{Cadre}
% ======================================================

\subsection{Conditionnement et sous-ensemble applicable}

Soit $n$ le nombre total de features binaires, et $F\subset\{1,\ldots,n\}$ l'ensemble
des $K=|F|$ indices gelés par la clause, chacun à une valeur $v_i^\star\in\{0,1\}$
($i\in F$) tirée au hasard à la création de la clause. On note
\[
  \DF = \bigl\{\, x\in\{0,1\}^n \;:\; x_i = v_i^\star \ \ \forall i\in F \,\bigr\}
\]
le \textbf{sous-ensemble applicable}. Dans l'implémentation, c'est exactement la
fonction \texttt{applies(x)} de \texttt{RandomStumpClause} : seuls les exemples de
$\DF$ sont vus par la clause, à l'entraînement comme à l'inférence.

Pour $i\in F$, les deux automates $v_i,\vbar_i$ sont gelés à l'état \Exc{} et exclus
de tout apprentissage (\texttt{active[i]=false}). Ils ne jouent aucun rôle dans la
suite : la feature $i$ est déjà « consommée » par le filtre \texttt{applies}.

On note $F^c = \{1,\ldots,n\}\setminus F$ l'ensemble des $n-K$ features
\textbf{libres}, chacune gouvernée par sa propre paire d'automates $(v_i,\vbar_i)_{i\in F^c}$,
mise à jour par la règle à 4 cas (Section~2 du document précédent), \textbf{restreinte
aux exemples de $\DF$} : l'entraînement tire ses exemples uniquement dans le
sous-ensemble applicable (\texttt{subsetIdx} dans le code).

\subsection{Hypothèse de résolution}

On suppose que le tirage des $K$ features conditionnées a réussi à isoler un
\textbf{unique littéral pertinent} : il existe $j\in F^c$ tel que, restreint à
$\DF$, le label $y$ ne dépend que de $x_j$, selon un canal de bruit
$a_j = P(y=1\mid x_j=1,\DF)$, $b_j = P(y=1\mid x_j=0,\DF)$, avec $a_j\ne b_j$
(cas IDENTITY si $a_j>b_j$, NOT si $b_j>a_j$), et $c_j = P(x_j=1\mid \DF)\in(0,1)$.

Pour toute autre feature libre $i\in F^c\setminus\{j\}$, on suppose qu'elle est
\textbf{non pertinente} restreinte à $\DF$ : conditionnellement indépendante de $y$,
c'est-à-dire
\[
  P(y=1\mid x_i=1,\DF) \;=\; P(y=1\mid x_i=0,\DF) \;=:\; p_y \in(0,1),
\]
avec une marginale propre $c_i = P(x_i=1\mid\DF)\in(0,1)$, pas nécessairement égale à
$\tfrac12$.

\begin{remark}
Cette hypothèse de résolution --- l'existence d'un $j$ isolé --- est exactement ce
que le conditionnement à $K$ features est censé produire : c'est un fait
\emph{statistique} sur le tirage aléatoire des $K$ conditions, pas une propriété
garantie a priori. Voir Section~\ref{sec:portee} pour la discussion complète.
\end{remark}

% ======================================================
\section{Indépendance des mises à jour entre littéraux libres}
% ======================================================

\begin{lemma}[Découplage multi-features]\label{lem:decouplage_multi}
Pour tout $i\in F^c$, l'action appliquée au couple $(v_i,\vbar_i)$ à l'instant $t$
ne dépend que de $(x_{i,t}, y_t)$ --- la valeur de la feature $i$ et le label de
l'exemple courant, tous deux restreints à $\DF$. Elle ne dépend :
\begin{itemize}
  \item ni de la sortie $C(x_t)$ de la clause,
  \item ni de l'état des automates $(v_{i'},\vbar_{i'})$ pour $i'\ne i$,
  \item ni des valeurs $x_{i'}$ pour $i'\ne i$.
\end{itemize}
\end{lemma}

\begin{proof}
Vérification directe sur \texttt{updateMasked} : la boucle de mise à jour traite
chaque feature libre $i$ indépendamment, en appliquant la table de la règle à 4 cas
(Section~2 du document précédent) au couple $(x_{i,t}, y_t)$ uniquement. Aucun terme
croisé entre features n'apparaît dans le code ni dans la spécification de la règle.
Le cas particulier $K=1$ (une seule feature libre observée) redonne exactement le
Lemme de découplage du document précédent.
\end{proof}

\begin{remark}[Ce que ce lemme ne suppose \emph{pas}]
On ne suppose \emph{aucune} indépendance statistique entre les $x_i$ eux-mêmes :
deux pixels voisins d'une image, par exemple, sont typiquement corrélés. Le
Lemme~\ref{lem:decouplage_multi} est un fait \emph{architectural} (comment la mise à
jour est câblée), pas une hypothèse sur les données. Ce qui compte pour la suite est
seulement la loi marginale du couple $(x_i, y)$ restreint à $\DF$ --- peu importe
ce qui se passe avec les autres $x_{i'}$.
\end{remark}

Conséquence directe : la chaîne $(v_i(t),\vbar_i(t))$ est, pour chaque $i\in F^c$
\emph{pris isolément}, exactement une instance de la chaîne bidirectionnelle étudiée
dans le document précédent (Lemme~4 : « chaîne bidirectionnelle », Lemme ergodique),
avec ses propres paramètres $(a_i,b_i,c_i)$ calculés sur $\DF$.

% ======================================================
\section{Convergence du littéral pertinent}
% ======================================================

\begin{theorem}[Rappel --- convergence du littéral pertinent]\label{thm:pertinent}
Sous l'hypothèse de résolution (Section~2.2), avec $\Delta_j =
(1-c_j)(1-b_j)-c_j(1-a_j)$ (cas IDENTITY, $a_j>b_j$, $\Delta_j>0$) et
\[
  S \;>\; S_{\text{bas}}(j) \;:=\; \frac{(1-c_j)\,b_j\,\alpha}{\Delta_j},
\]
alors $\rho_{v_j}>1$ et $\rho_{\vbar_j}<1$, et
\[
  \pi_{v_j}(\Inc) \xrightarrow[N\to\infty]{} 1, \qquad
  \pi_{\vbar_j}(\Exc) \xrightarrow[N\to\infty]{} 1.
\]
(Cas NOT : symétrique, avec $\Delta_j'=-\Delta_j$ et le seuil correspondant.)
\end{theorem}

\begin{proof}
Application directe du Théorème « Convergence IDENTITY, cas bruité » (resp. « NOT »)
du document précédent, avec $c,a,b,\Delta$ remplacés par $c_j,a_j,b_j,\Delta_j$
--- les statistiques locales de $(x_j,y)$ restreintes à $\DF$. La preuve est
identique terme à terme : elle ne repose que sur la loi marginale de $(x_j,y)$, qui
est ici la loi \emph{conditionnelle} à $\DF$ plutôt que la loi globale, ce qui ne
change rien à l'argument (bilan détaillé, chaîne ergodique).
\end{proof}

% ======================================================
\section{Neutralité des littéraux non pertinents}
% ======================================================

C'est le résultat nouveau de ce document. Il répond à une question posée en
pratique (par Florent Avellaneda) : \emph{est-ce que les probabilités $S,\alpha$
de la règle peuvent faire converger un littéral non pertinent vers une contrainte
fausse, simplement à cause d'un déséquilibre marginal de sa feature ?}

\subsection{Calcul des rapports $\rho$}

Pour $i\in F^c\setminus\{j\}$, non pertinente ($a_i=b_i=p_y$), les probabilités de
transition (Section~5 du document précédent, avec $a=b=p_y$) deviennent :
\begin{align*}
  p^+_{v_i} &= (1-c_i)(1-p_y)\,S, &
  p^-_{v_i} &= (1-c_i)\,p_y\,\alpha + c_i(1-p_y)\,S; \\
  p^+_{\vbar_i} &= c_i(1-p_y)\,S, &
  p^-_{\vbar_i} &= (1-c_i)(1-p_y)\,S + c_i\,p_y\,\alpha.
\end{align*}

\begin{lemma}[Signe de $\rho_{v_i}-1$ et $\rho_{\vbar_i}-1$]\label{lem:signe_neutre}
\[
  \rho_{v_i} < 1
  \;\Longleftrightarrow\;
  (1-2c_i)(1-p_y)\,S \;<\; (1-c_i)\,p_y\,\alpha,
  \qquad
  \rho_{\vbar_i} < 1
  \;\Longleftrightarrow\;
  (2c_i-1)(1-p_y)\,S \;<\; c_i\,p_y\,\alpha.
\]
\end{lemma}

\begin{proof}
$\rho_{v_i}<1 \iff p^+_{v_i}<p^-_{v_i}$ (car $p^-_{v_i}>0$ toujours) :
\[
  (1-c_i)(1-p_y)S < (1-c_i)p_y\alpha + c_i(1-p_y)S
  \;\Longleftrightarrow\;
  (1-p_y)S\bigl[(1-c_i)-c_i\bigr] < (1-c_i)p_y\alpha,
\]
d'où la première inégalité, le facteur $(1-c_i)-c_i=1-2c_i$. Calcul symétrique pour
$\rho_{\vbar_i}<1$, en factorisant $c_i-(1-c_i)=2c_i-1$.
\end{proof}

\begin{theorem}[Neutralité --- fenêtre supérieure sur $S$]\label{thm:neutralite}
Soit $i\in F^c\setminus\{j\}$ non pertinente, de marginale $c_i\in(0,1)$. Posons
\[
  S_i^{\max} \;:=\;
  \begin{cases}
    \dfrac{c_i\,p_y\,\alpha}{(2c_i-1)(1-p_y)} & \text{si } c_i > \tfrac12, \\[10pt]
    \dfrac{(1-c_i)\,p_y\,\alpha}{(1-2c_i)(1-p_y)} & \text{si } c_i < \tfrac12, \\[10pt]
    +\infty & \text{si } c_i = \tfrac12.
  \end{cases}
\]
Si $S < S_i^{\max}$, alors $\rho_{v_i}<1$ \emph{et} $\rho_{\vbar_i}<1$
simultanément, donc
\[
  \pi_{v_i}(\Exc) \xrightarrow[N\to\infty]{} 1
  \qquad\text{et}\qquad
  \pi_{\vbar_i}(\Exc) \xrightarrow[N\to\infty]{} 1,
\]
c'est-à-dire que le littéral $i$ converge vers la configuration $(\Exc,\Exc)$
--- \textbf{tautologique}, ne contraignant rien.
\end{theorem}

\begin{proof}
\textbf{Cas $c_i=\tfrac12$.} Les deux membres de gauche du
Lemme~\ref{lem:signe_neutre} s'annulent ($1-2c_i=2c_i-1=0$), et les deux membres de
droite sont strictement positifs (car $p_y\in(0,1)$, $\alpha>0$, $c_i,1-c_i>0$).
Les deux inégalités sont donc vraies pour \emph{tout} $S>0$ : $\rho_{v_i}<1$ et
$\rho_{\vbar_i}<1$ inconditionnellement.

\textbf{Cas $c_i>\tfrac12$.} Le facteur $1-2c_i<0$ rend la première inégalité du
Lemme~\ref{lem:signe_neutre} automatiquement vraie (membre de gauche négatif,
membre de droite positif) : $\rho_{v_i}<1$ pour tout $S$. Pour la seconde, le
facteur $2c_i-1>0$ impose de diviser --- sans changer le sens de l'inégalité ---
pour isoler $S$ :
\[
  (2c_i-1)(1-p_y)\,S < c_i\,p_y\,\alpha
  \;\Longleftrightarrow\;
  S < \frac{c_i\,p_y\,\alpha}{(2c_i-1)(1-p_y)} = S_i^{\max}.
\]

\textbf{Cas $c_i<\tfrac12$.} Symétrique (échanger les rôles de $v_i$ et $\vbar_i$),
la seconde inégalité est automatique et la première donne le seuil correspondant.

Dans les trois cas, $S<S_i^{\max}$ entraîne $\rho_{v_i}<1$ et $\rho_{\vbar_i}<1$
simultanément. Le Lemme « chaîne bidirectionnelle » du document précédent
s'applique alors aux deux chaînes (chacune ergodique, espace d'états fini), donnant
$\pi_{v_i}(\Exc) = 1/(\rho_{v_i}^{-N}+1)$ et de même pour $\vbar_i$, toutes deux
$\to 1$ quand $N\to\infty$.
\end{proof}

\begin{remark}[Interprétation --- le bruit vient du déséquilibre, pas de la
corrélation]
Le résultat est contre-intuitif au premier abord : la feature $i$ est \emph{parfaitement
indépendante} de $y$, et pourtant sa convergence vers la neutralité n'est pas
automatique. La cause n'est pas une corrélation cachée avec $y$ : c'est le simple
\textbf{déséquilibre marginal} $c_i\ne\tfrac12$ (la feature vaut plus souvent 1 que
0, ou l'inverse) qui, combiné à un $S$ trop grand, fait dériver un des deux
automates vers \Inc{} --- une inclusion \emph{fantôme}, fondée sur rien d'autre que
la fréquence brute de la feature. C'est exactement la question posée par Florent :
plus il y a de features libres, plus il y a de chances qu'au moins une ait un
déséquilibre marginal significatif, et donc plus la fenêtre tolérable pour $S$ se
resserre.
\end{remark}

% ======================================================
\section{Convergence globale de la clause conditionnée}
% ======================================================

\begin{theorem}[Convergence de la clause à $K$ features]\label{thm:global}
Sous l'hypothèse de résolution (Section~2.2), supposons que $S$ satisfasse
simultanément
\[
  \boxed{\;S_{\text{bas}}(j) \;<\; S \;<\; \min_{i\in F^c\setminus\{j\}} S_i^{\max}\;}
\]
où $S_{\text{bas}}(j)$ est le seuil du Théorème~\ref{thm:pertinent} et $S_i^{\max}$
celui du Théorème~\ref{thm:neutralite}. Alors, quand $N\to\infty$ :
\begin{itemize}
  \item $(v_j,\vbar_j)$ converge vers la configuration codant \Inc/\Exc{} ou
  \Exc/\Inc{} selon IDENTITY/NOT (Théorème~\ref{thm:pertinent}) ;
  \item pour tout $i\in F^c\setminus\{j\}$, $(v_i,\vbar_i)$ converge vers
  $(\Exc,\Exc)$ (Théorème~\ref{thm:neutralite}) ;
\end{itemize}
simultanément, avec probabilité stationnaire jointe $\to 1$. La clause, restreinte
à $\DF$, converge donc vers l'implémentation correcte du littéral unique $x_j$ (ou
$\neg x_j$), sans contrainte parasite issue des $n-K-1$ autres features libres.
\end{theorem}

\begin{proof}
Notons $q_i(N) = 1-\pi_i(\text{configuration correcte pour }i)$ l'erreur
stationnaire de la feature libre $i$ --- où « correcte » signifie \Inc/\Exc{}
adéquat pour $i=j$ (Théorème~\ref{thm:pertinent}), ou $(\Exc,\Exc)$ pour
$i\ne j$ (Théorème~\ref{thm:neutralite}). Chacune des $n-K$ chaînes
$(v_i,\vbar_i)_{i\in F^c}$ est ergodique (Lemme « chaîne bidirectionnelle », état
fini, irréductible et apériodique), donc admet une unique loi stationnaire, et la
borne géométrique du document précédent donne $q_i(N) \le \rho_i^{-N}$ pour un
$\rho_i>1$ approprié (le $\rho_v$ ou $\rho_{\vbar}$ pertinent selon le cas), avec
$\rho_i>1$ garanti par les Théorèmes~\ref{thm:pertinent} et \ref{thm:neutralite}
sous la condition encadrée sur $S$.

Il y a $n-K$ features libres, donc $n-K$ termes $q_i(N)$ (un seul pour $j$, deux
--- $v_i$ et $\vbar_i$ --- pour chaque $i\ne j$, soit au plus $2(n-K)-1$ chaînes
au total, un nombre \emph{fini} fixé par $n$ et $K$). Par l'inégalité de
Boole (union finie) :
\[
  P(\text{au moins une chaîne dans la mauvaise configuration})
  \;\le\; \sum_{\text{chaînes}} q_i(N)
  \;\xrightarrow[N\to\infty]{}\; 0,
\]
puisque c'est une somme \emph{finie} de termes tendant individuellement vers $0$.
Donc
\[
  P(\text{toutes les chaînes dans la configuration correcte}) \xrightarrow[N\to\infty]{} 1,
\]
ce qui est exactement l'énoncé du théorème.
\end{proof}

\begin{remark}[Pourquoi une intersection finie suffit]
Le nombre de features libres $n-K$ est fixé dès la création de la clause --- ce
n'est pas une limite $n\to\infty$. L'union bound ci-dessus est donc une somme
finie de termes qui décroissent chacun géométriquement en $N$ ($\rho_i^{-N}$) :
la convergence jointe est \emph{elle-même} géométrique en $N$, avec un taux
gouverné par le pire des $\rho_i$ (le plus proche de $1$ parmi le littéral
pertinent et tous les littéraux neutres). C'est une conséquence directe et
quantifiée du Lemme~\ref{lem:decouplage_multi} : parce que les mises à jour sont
découplées entre littéraux, la convergence jointe ne coûte rien de plus qu'une
somme finie d'erreurs individuelles.
\end{remark}

\begin{corollary}[Existence d'une fenêtre non vide]
La fenêtre $\bigl(S_{\text{bas}}(j),\,\min_i S_i^{\max}\bigr)$ du
Théorème~\ref{thm:global} est non vide dès que
$S_{\text{bas}}(j) < \min_i S_i^{\max}$, une condition qui dépend uniquement des
statistiques locales $(a_j,b_j,c_j)$ et $(c_i)_{i\ne j}$ sur $\DF$ --- calculable
explicitement à partir des données, exactement comme le seuil $S^\star$ du
Test~A du document précédent. Rien ne garantit que cette fenêtre soit non vide
pour tout jeu de données : si une feature non pertinente a un déséquilibre
marginal $c_i$ très fort, $S_i^{\max}$ peut devenir trop petit pour laisser de la
place au-dessus de $S_{\text{bas}}(j)$ --- auquel cas \emph{aucun} choix de $S$ ne
fait converger la clause correctement, et augmenter $\alpha$ (qui apparaît au
numérateur de $S_i^{\max}$ et de $S_{\text{bas}}(j)$ de façon symétrique) ne
suffit pas non plus à rouvrir la fenêtre à lui seul.
\end{corollary}

% ======================================================
\section{Validation empirique}\label{sec:validation}
% ======================================================

Les deux théorèmes centraux (\ref{thm:pertinent} et \ref{thm:neutralite}) ont
depuis été confrontés à des données réelles, avec un accord net.

\subsection{La fenêtre sur $S$, mesurée sur XOR}

Sur \texttt{NoisyXORTrainingData.txt} ($K=1$, feature conditionnée $=x_1$,
littéral pertinent $=x_2$), la mesure directe donne $c_j=0{,}5$, $a_j=0{,}603$,
$b_j=0{,}397$ (bruit $\approx 39{,}7\,\%$). Le Théorème~\ref{thm:pertinent}
prédit un seuil $S_{\text{bas}}(j)$ qui, réécrit comme seuil sur $\alpha$ à $S$
fixé, donne $\alpha^\star(S)=0{,}519\,S$. Un balayage de 42 configurations
$(S,\alpha)$, 20 runs chacune, confirme cette prédiction à la configuration
près : l'accuracy s'effondre exactement (et de façon reproductible, variance
nulle sur 20 runs) pour $\alpha>\alpha^\star(S)$, et reste correcte en-deçà,
pour les 6 valeurs de $S$ testées. Voir
\textit{Rôle de $S$ et $\alpha$ : théorie et validation} pour le détail complet
(tableau, calculs, comparaison configuration par configuration).

\subsection{La neutralité, mesurée directement (Tests D/E)}

Le Théorème~\ref{thm:neutralite} a été testé isolément (\texttt{n=2} features,
une pertinente et une non pertinente de marginale $c_i$ contrôlée) dans
\texttt{validation\_theorie\_neutralite.cpp}. Seuil prédit $S_1^{\max}=0{,}20$
(Test D, $c_1=0{,}8$) : bascule empirique mesurée entre $S=0{,}18$ et $S=0{,}22$.
Seuil prédit $c_1^\star=0{,}588$ (Test E, $S=0{,}5$ fixé) : bascule mesurée
entre $c_1=0{,}55$ et $c_1=0{,}65$. Dans les deux cas, la zone de transition
empirique encadre exactement la prédiction théorique.

\subsection{Pourquoi Iris et Digits ne montrent presque aucun effet}

Le Corollaire (fenêtre non vide) prévient que $S_{\text{bas}}(j)$ peut être
très inférieur à $\min_i S_i^{\max}$ selon les données --- rendant la fenêtre
si large qu'aucune valeur de $S,\alpha\in(0,1]$ ne s'en approche. C'est
exactement ce qui est mesuré sur Iris et Digits : contrairement à XOR (où
$a_j-b_j=0{,}206$, un signal faible), plusieurs features d'Iris ont un écart
$a-b$ très supérieur (jusqu'à $0{,}87$ pour une feature de Digits), donnant des
$\alpha^\star(S{=}1)$ de l'ordre de $2$ à $230$ --- très au-delà de $1$. Un
balayage complet sur ces deux datasets confirme : la destination ne bascule
quasiment jamais dans $(0,1]$, et la légère dégradation observée à petit
$S$/grand $\alpha$ est un effet de \emph{vitesse} (Section~6.6 du document
précédent), pas de destination --- cohérent avec le Théorème~\ref{thm:pertinent}
qui ne garantit la destination qu'\emph{asymptotiquement en $N$}, pas
instantanément à budget d'entraînement fini.

\subsection{Confirmation indépendante : $K=0$ échoue là où l'hypothèse de résolution ne peut pas s'appliquer}

Section~\ref{sec:portee} note que ce document ne couvre pas les cibles à
$m\ge2$ littéraux sans conditionnement. Un test direct ($K=0$ sur Iris, aucune
feature gelée) confirme cette limite de façon indépendante :
\textbf{22,7\,\% de précision, sous le niveau du hasard (33\,\% pour 3
classes)}. Ce n'est pas dû à un manque de signal --- plusieurs features d'Iris
sont individuellement très informatives (Section~\ref{sec:validation},
ci-dessus) --- mais au fait que distinguer une espèce nécessite de combiner
\emph{plusieurs} littéraux, une cible hors de portée de la règle à 4 cas sans
le conditionnement qui la réduit à l'hypothèse de résolution
(Section~\ref{sec:portee}, ci-dessous). C'est la seconde confirmation
indépendante de cette limite, après \texttt{hybrid\_coord\_4case.cpp}.

% ======================================================
\section{Portée et limites}\label{sec:portee}
% ======================================================

\subsection{Ce qui reste hors de la preuve : le tirage du conditionnement}

Le Théorème~\ref{thm:global} suppose l'hypothèse de résolution : qu'un tirage
aléatoire particulier des $K$ features et valeurs conditionnées a isolé un unique
littéral pertinent $j$. Rien dans ce document ne prouve qu'un tirage aléatoire
\emph{trouve} un tel $j$ --- c'est un phénomène combinatoire distinct, caractérisé
empiriquement dans ce travail sous le nom de \emph{loterie combinatoire} :
\begin{itemize}
  \item Mux11 ($K=3$, $\binom{11}{3}=165$ combinaisons, $1$ utile) : $M=300$
  insuffisant (89\,\%), $M=8000$ résout (99,97\,\%).
  \item MNIST ($K=2$, $\binom{784}{2}\approx 307\,000$ combinaisons) : amélioration
  continue mais lente avec $M$ (85,7\,\% à $M{=}500$, 94,9\,\% à $M{=}8000$),
  cohérente avec un espace de combinaisons bien plus grand que Mux11 ou Digits.
\end{itemize}
Une preuve complète du pipeline demanderait de combiner ce document (convergence
\emph{une fois} le conditionnement fixé) avec un argument probabiliste séparé
(type coupon collector / loi binomiale) bornant $P(\text{au moins une clause parmi
$M$ tire un conditionnement résolvant})$ --- non traité ici.

\subsection{Ce qui reste hors de la preuve : les cibles à $m\ge2$ littéraux sans conditionnement}

Ce document ne prouve --- et ne cherche pas à prouver --- la convergence directe
vers une cible à $m\ge2$ littéraux sans passer par un conditionnement qui la
réduit d'abord à un littéral unique sur $\DF$. C'est précisément le régime que
couvre PCL (Belaid et al., AAAI 2025) par un argument différent (probabilité
d'inclusion unique $p_k$, sans conditionnement ni routeur). L'échec empirique de
\texttt{hybrid\_coord\_4case.cpp} (plafond à $\sim$51\,\% sur une cible à 2
littéraux, sans conditionnement) suggère que l'asymétrie propre de la règle à 4
cas (Lemme de découplage : $y=1$ ne pousse jamais activement vers \Inc) rend ce
régime hors de portée pour notre règle spécifique --- notre stratégie est
structurellement différente : réduire la cible par conditionnement plutôt que
la résoudre directement.

\subsection{Ce qui est acquis}

En comparaison de PCL, ce document fournit, dans son régime propre (conditionnement
$+$ littéral unique résiduel) :
\begin{itemize}
  \item une fenêtre \emph{explicite et calculable} $(S_{\text{bas}}(j), \min_i
  S_i^{\max})$ sur l'hyperparamètre $S$, avec formule fermée des deux bornes ;
  \item un \emph{taux de convergence} explicite (géométrique en $N$, gouverné par
  $\min(\rho_j, \min_i \rho_i)$) ;
  \item une explication quantifiée du lien entre nombre de features libres et
  « bruit perçu » (Théorème~\ref{thm:neutralite}) --- question posée
  indépendamment par F.~Avellaneda, et qui se trouve avoir une réponse fermée dans
  ce cadre.
\end{itemize}

\end{document}
